% !TEX program = pdflatex
\documentclass[11pt,a4paper]{article}

% =========================================================================
% INFORME DE CONTROL DEL PROYECTO
% Proyecto: Horus Sentinel
% Gestión de Proyectos TICS (EST315) -- Semestre 2026-II -- Team Sobrevivientes
% =========================================================================

% --- PAQUETES DE CONFIGURACIÓN ---
\usepackage[utf8]{inputenc}
\usepackage[T1]{fontenc}
\usepackage[spanish,es-tabla]{babel}
% Compatibilidad con babel-spanish reciente (flechas de TikZ y signo de porcentaje)
\AtBeginDocument{\shorthandoff{<>}\spanishplainpercent}
\usepackage{geometry}
\geometry{
    top=3.6cm,
    bottom=2.3cm,
    left=2.2cm,
    right=2.2cm,
    % PARTE MODIFICABLE: headheight debe ser mayor que la altura del logo más alto (2.56cm = 73pt).
    headheight=80pt,
    headsep=15pt,
    footskip=20pt
}

\usepackage{graphicx}
\usepackage{fancyhdr}
\usepackage[table]{xcolor}
\usepackage{titlesec}
\usepackage{booktabs}
\usepackage{tabularx}
\usepackage{array}
\usepackage{amsmath}
\usepackage{enumitem}
\usepackage{microtype}
\usepackage{amssymb}
% --- RECURSOS VISUALES: GRÁFICOS Y CÓDIGO DE COLORES ---
\usepackage{tikz}
\usetikzlibrary{positioning,calc}
\definecolor{GrisFila}{RGB}{246, 248, 251}
\definecolor{CuadCerca}{RGB}{150, 30, 30}
\definecolor{CuadSatisf}{RGB}{180, 130, 20}
\definecolor{CuadInfo}{RGB}{10, 45, 90}
\definecolor{CuadMonit}{RGB}{110, 115, 125}
\newcommand{\etiqueta}[2]{{\setlength{\fboxsep}{2pt}\colorbox{#1}{\textcolor{white}{\scriptsize\bfseries #2}}}}
\newcommand{\cerca}{\etiqueta{CuadCerca}{Gestionar de cerca}}
\newcommand{\satisf}{\etiqueta{CuadSatisf}{Mantener satisfecho}}
\newcommand{\informado}{\etiqueta{CuadInfo}{Mantener informado}}
\newcommand{\monit}{\etiqueta{CuadMonit}{Monitorear}}
\usepackage{needspace}
\usepackage{longtable}
\definecolor{Rama1}{RGB}{10, 45, 90}
\definecolor{Rama2}{RGB}{40, 95, 150}
\definecolor{Rama3}{RGB}{180, 130, 20}
\definecolor{Rama4}{RGB}{30, 110, 60}
\definecolor{Rama5}{RGB}{110, 60, 140}
\definecolor{RaciR}{RGB}{150, 30, 30}
\definecolor{RaciA}{RGB}{10, 45, 90}
\definecolor{RaciC}{RGB}{180, 130, 20}
\definecolor{RaciI}{RGB}{140, 145, 155}
\definecolor{NivAlto}{RGB}{30, 110, 60}
\definecolor{NivMedio}{RGB}{180, 130, 20}
\definecolor{NivBajo}{RGB}{150, 30, 30}
\usetikzlibrary{patterns}
\definecolor{SevBaja}{RGB}{30, 110, 60}
\definecolor{SevMedia}{RGB}{180, 130, 20}
\definecolor{SevAlta}{RGB}{150, 30, 30}
\usetikzlibrary{arrows.meta,shapes.geometric,fit,backgrounds}
\usepackage{listings}
\lstset{basicstyle=\ttfamily\scriptsize, frame=single, rulecolor=\color{AzulUNA!40}, backgroundcolor=\color{GrisFila}, columns=fullflexible, keepspaces=true, breaklines=true, xleftmargin=4pt, xrightmargin=4pt, literate={á}{{\'a}}1 {é}{{\'e}}1 {í}{{\'i}}1 {ó}{{\'o}}1 {ú}{{\'u}}1 {ñ}{{\~n}}1 {Á}{{\'A}}1}

\usepackage{hyperref}

% Logos: funciona a 2 o 3 niveles de profundidad dentro de docs/latex/
\graphicspath{{../assets/img/}{../../assets/img/}{../../../assets/img/}}

% --- PALETA INSTITUCIONAL FIEI / UNA PUNO ---
\definecolor{AzulUNA}{RGB}{10, 45, 90}
\definecolor{DoradoFIEI}{RGB}{180, 130, 20}
\definecolor{GrisOscuro}{RGB}{50, 50, 50}
\definecolor{AzulClaro}{RGB}{228, 236, 246}

% --- CONFIGURACIÓN DE LISTAS ---
\setlist{topsep=2pt, itemsep=2pt, parsep=1pt}

% --- COLUMNAS DE TABLA ---
\newcolumntype{L}{>{\raggedright\arraybackslash}X}
\newcolumntype{P}[1]{>{\raggedright\arraybackslash}p{#1}}
\newcolumntype{C}[1]{>{\centering\arraybackslash}p{#1}}
\newcolumntype{Q}[1]{>{\raggedleft\arraybackslash}p{#1}}
\newcolumntype{R}{>{\raggedleft\arraybackslash}X}
\renewcommand{\arraystretch}{1.25}

% --- FORMATO DE ENLACES ---
\hypersetup{
    colorlinks=true,
    linkcolor=AzulUNA,
    urlcolor=AzulUNA,
    citecolor=AzulUNA,
    pdftitle={Informe De Control Del Proyecto -- Horus Sentinel},
    pdfauthor={Team Sobrevivientes}
}

% --- FORMATO DE TÍTULOS ---
\titleformat{\section}{\color{AzulUNA}\normalfont\large\bfseries}{\thesection.}{0.7em}{}[\color{AzulUNA}\titlerule]
\titleformat{\subsection}{\color{AzulUNA}\normalfont\normalsize\bfseries}{\thesubsection.}{0.5em}{}
\titlespacing*{\section}{0pt}{9pt}{4pt}
\titlespacing*{\subsection}{0pt}{6pt}{2pt}

% --- ENCABEZADO Y PIE DE PÁGINA (fancyhdr) ---
\pagestyle{fancy}
\fancyhf{}

% ===== TAMAÑO DE LOS LOGOS (PARTE MODIFICABLE) =====
% logo_fiei.png  : height=1.45cm (lado izquierdo)
% logo_equipo.png: height=2.56cm (lado derecho; 25 % más grande porque la imagen tiene más borde blanco)
% Para cambiar el tamaño, modifique solo el valor height=... de cada \includegraphics.
% Si el logo más alto supera 2.6cm, aumente también headheight y top en la configuración de página.
\fancyhead[L]{%
    \includegraphics[height=1.45cm,keepaspectratio]{logo_fiei.png}%
}
\fancyhead[C]{%
    \scriptsize \textbf{\textcolor{AzulUNA}{UNIVERSIDAD NACIONAL DEL ALTIPLANO DE PUNO}}\\[2pt]
    \tiny \textcolor{GrisOscuro}{FACULTAD DE INGENIERÍA ESTADÍSTICA E INFORMÁTICA}\\[2pt]
    \scriptsize \textbf{\textcolor{AzulUNA}{Gestión de Proyectos TICS (EST315) \textbar\ Team Sobrevivientes}}\\[2pt]
    \footnotesize \textbf{\textcolor{DoradoFIEI}{INFORME DE CONTROL DEL PROYECTO}}%
}
\fancyhead[R]{%
    \includegraphics[height=2.56cm,keepaspectratio]{logo_equipo.png}%
}

\renewcommand{\headrulewidth}{0pt}
\renewcommand{\headrule}{}

\fancyfoot[L]{\scriptsize \textcolor{GrisOscuro}{Horus Sentinel}}
\fancyfoot[C]{\small \thepage}
\fancyfoot[R]{\scriptsize \textcolor{GrisOscuro}{Puno, Perú -- Semestre 2026-II}}
\renewcommand{\footrulewidth}{0.4pt}

\begin{document}

\begin{center}
    {\Large \textbf{\textcolor{AzulUNA}{INFORME DE CONTROL DEL PROYECTO}}}\\[3pt]
    {\normalsize \textbf{Horus Sentinel}}\\[4pt]
    \footnotesize \textbf{Fecha:} lunes 21/09/2026 (semana 5) \quad \textbar\quad \textbf{Corte:} viernes 18/09/2026 (fin de la semana 4) \quad \textbar\quad \textbf{Fase:} Monitoreo
\end{center}

\vspace{0.1cm}

\section{Datos Generales}
\noindent\begin{tabularx}{\textwidth}{@{}P{4.2cm} L@{}}
\toprule
\textbf{Proyecto:} & Horus Sentinel \\
\textbf{Encargado por:} & Clyde Neil Paricahua Pari -- Project Manager, en la revisión del sprint del viernes 18/09/2026 \\
\textbf{Elaborado por:} & Ruth Karina Apaza Solis (valor ganado y horas), Eliseo Tarqui Ajahuana (estado técnico) y Viviana Mary Quisbert Quispe (riesgos) \\
\textbf{Insumos:} & Líneas base del cronograma y de costos, registro de horas en los Issues, gastos con comprobante y registro de riesgos \\
\textbf{Destinatario:} & El equipo como sponsor, en la reunión de control del martes 22/09/2026 \\
\bottomrule
\end{tabularx}

\section{Encargo del Project Manager}
Al aprobarse el cronograma y el presupuesto detallados, el Project Manager pidió medir el estado real de Horus Sentinel antes de comprometer más horas del equipo:

\vspace{0.1cm}
{\small
\rowcolors{2}{white}{GrisFila}
\noindent\begin{tabularx}{\textwidth}{@{}C{0.7cm} L P{3.4cm} C{1.9cm}@{}}
\toprule
\rowcolor{AzulClaro}
\textbf{N.\textsuperscript{o}} & \textbf{Tarea encargada} & \textbf{Responsable} & \textbf{Plazo} \\
\midrule
1 & Calcular el valor ganado al corte de la semana 4 con el avance real de cada actividad. & Ruth K. Apaza Solis & Dom 20/09 \\
2 & Comparar las horas registradas con las horas planificadas. & Ruth K. Apaza Solis & Dom 20/09 \\
3 & Informar el avance de la capacitación geoespacial y probar el cálculo de rutas en el servidor. & Eliseo Tarqui Ajahuana & Dom 20/09 \\
4 & Actualizar el registro de riesgos con lo ocurrido en las semanas 2 a 4. & Viviana M. Quisbert Quispe & Dom 20/09 \\
5 & Comparar la rentabilidad de Horus Sentinel con la de CasaScore con las mismas condiciones. & Ruth K. Apaza Solis & Lun 21/09 \\
\bottomrule
\end{tabularx}
}

\Needspace{20\baselineskip}
\section{Valor Ganado al Corte}
Valores en soles. El presupuesto del trabajo (BAC) es S/ 11,829.93; la contingencia se controla aparte. Los costos reales son simulados a partir del registro de horas y gastos.

\vspace{0.1cm}
{\footnotesize
\setlength{\tabcolsep}{3pt}
\rowcolors{2}{white}{GrisFila}
\noindent\begin{tabularx}{\textwidth}{@{}C{0.7cm} L Q{1.3cm} Q{1.3cm} C{0.9cm} Q{1.3cm} Q{1.3cm} Q{1.3cm} Q{1.15cm}@{}}
\toprule
\rowcolor{AzulClaro}
\textbf{Cód.} & \textbf{Actividad} & \textbf{BAC} & \textbf{PV} & \textbf{\% real} & \textbf{EV} & \textbf{AC} & \textbf{SV} & \textbf{CV} \\
\midrule
2.1 & Ficha técnica y requerimientos & 372.65 & 372.65 & 100\,\% & 372.65 & 372.65 & 0.00 & 0.00 \\
1.1 & Viabilidad y Acta de Constitución & 372.65 & 372.65 & 100\,\% & 372.65 & 380.10 & 0.00 & \textcolor{SevAlta}{-7.45} \\
2.3 & Capacitación geoespacial & 372.65 & 372.65 & 55\,\% & 204.96 & 241.10 & \textcolor{SevAlta}{-167.69} & \textcolor{SevAlta}{-36.14} \\
1.2 & Planificación base & 496.86 & 488.31 & 85\,\% & 422.33 & 455.20 & \textcolor{SevAlta}{-65.98} & \textcolor{SevAlta}{-32.87} \\
2.2 & Modelo de datos geoespacial & 496.86 & 186.00 & 10\,\% & 49.69 & 58.30 & \textcolor{SevAlta}{-136.31} & \textcolor{SevAlta}{-8.61} \\
\midrule
\rowcolor{white}\multicolumn{3}{@{}l}{\textbf{Total al corte}} & \textbf{1,792.26} & & \textbf{1,422.28} & \textbf{1,507.35} & \textbf{\textcolor{SevAlta}{-369.98}} & \textbf{\textcolor{SevAlta}{-85.07}} \\
\bottomrule
\end{tabularx}
}

\Needspace{18\baselineskip}
\section{Indicadores y Pronóstico}
{\small
\rowcolors{2}{white}{GrisFila}
\noindent\begin{tabularx}{\textwidth}{@{}P{4.6cm} L Q{2.3cm} C{1.6cm}@{}}
\toprule
\rowcolor{AzulClaro}
\textbf{Indicador} & \textbf{Cálculo} & \textbf{Resultado} & \textbf{Estado} \\
\midrule
Índice del cronograma (SPI) & 1,422.28 / 1,792.26 & 0.79 & \etiqueta{SevAlta}{ROJO} \\
Índice del costo (CPI) & 1,422.28 / 1,507.35 & 0.94 & \etiqueta{SevMedia}{ÁMBAR} \\
Costo estimado al terminar (EAC) & 11,829.93 / 0.94 & S/ 12,537.51 & \etiqueta{SevAlta}{ROJO} \\
Variación al terminar (VAC) & 11,829.93 $-$ 12,537.51 & S/ $-$707.58 & -- \\
EAC si sigue el mismo ritmo & 1,507.35 + 10,407.65 / (0.94 $\times$ 0.79) & S/ 15,406.80 & \etiqueta{SevAlta}{ROJO} \\
Duración pronosticada & 18.75 semanas / 0.79 & 23.6 semanas & \etiqueta{SevAlta}{ROJO} \\
Fecha de término pronosticada & Inicio 31/08/2026 + 23.6 semanas & 12/02/2027 & -- \\
\bottomrule
\end{tabularx}
}

\vspace{0.1cm}
\noindent El EAC de S/ 12,537.51 supera la línea base de costos (S/ 12,421.43) aun usando toda la contingencia. Si el ritmo actual se mantiene, el costo llegaría a S/ 15,406.80, por encima del presupuesto total con reserva de gestión (S/ 13,663.57). El proyecto terminaría siete semanas después del cierre del semestre (24/12/2026).

\Needspace{14\baselineskip}
\section{Horas del Equipo}
{\small
\rowcolors{2}{white}{GrisFila}
\noindent\begin{tabularx}{\textwidth}{@{}L Q{2.5cm} Q{2.5cm} Q{2.5cm}@{}}
\toprule
\rowcolor{AzulClaro}
\textbf{Periodo} & \textbf{Planificadas} & \textbf{Registradas} & \textbf{Cumplimiento} \\
\midrule
Semanas 2 a 4 (31/08 -- 18/09) & 115 h & 97 h & 84\,\% \\
\bottomrule
\end{tabularx}
}

\vspace{0.1cm}
\begin{itemize}[leftmargin=*]
    \item Cuando ambos proyectos compitieron por el tiempo, el equipo atendió primero a CasaScore, como lo establece la evaluación de cartera.
    \item Las horas registradas produjeron menos avance del previsto: la capacitación geoespacial tomó más horas de las estimadas y retrasó el modelo de datos, que es la siguiente actividad de la ruta crítica.
    \item Sostener los dos proyectos exige 16 horas semanales por integrante; en las semanas 2 a 4 el promedio real fue menor, sin contar las demás asignaturas.
\end{itemize}

\Needspace{12\baselineskip}
\section{Estado Técnico}
{\small
\rowcolors{2}{white}{GrisFila}
\noindent\begin{tabularx}{\textwidth}{@{}P{3.4cm} L C{1.9cm}@{}}
\toprule
\rowcolor{AzulClaro}
\textbf{Aspecto} & \textbf{Situación al 18/09} & \textbf{Estado} \\
\midrule
Capacitación geoespacial & Se completaron los tutoriales de mapas y grafos de calles; el mapa de calor (KDE) aún no se ejecuta con datos de prueba. & \etiqueta{SevAlta}{Atrasado} \\
Servidor & El cálculo de rutas con el grafo del cercado de Puno no funciona en un plan gratuito; requiere el servidor de 8 GB (US\$ 40 al mes) desde octubre. & \etiqueta{SevMedia}{Confirmado} \\
Datos & No hay datos reales de incidentes; los reportes ciudadanos llegarían recién con la app terminada. & \etiqueta{SevMedia}{Pendiente} \\
Convenio municipal & No se ha iniciado el contacto con la municipalidad; el trámite tomaría al menos 3 semanas. & \etiqueta{SevAlta}{Sin iniciar} \\
\bottomrule
\end{tabularx}
}

\Needspace{14\baselineskip}
\section{Riesgos Actualizados}
{\small
\rowcolors{2}{white}{GrisFila}
\noindent\begin{tabularx}{\textwidth}{@{}L C{1.3cm} C{1.3cm} P{3.3cm}@{}}
\toprule
\rowcolor{AzulClaro}
\textbf{Riesgo} & \textbf{Antes} & \textbf{Ahora} & \textbf{Situación} \\
\midrule
El cronograma no alcanza y exige más jornadas & 60\,\% & 100\,\% & Ocurrido: SPI de 0.79 \\
El servidor necesita más memoria & 40\,\% & 40\,\% & Se mantiene \\
Pocos reportes ciudadanos & 40\,\% & 60\,\% & Sin campaña ni app \\
El convenio municipal no se firma & -- & 50\,\% & Nuevo riesgo \\
Alza de precios de alimentación y transporte & 30\,\% & 30\,\% & Se mantiene \\
\bottomrule
\end{tabularx}
}

\Needspace{16\baselineskip}
\section{Alternativas Evaluadas}
{\small
\rowcolors{2}{white}{GrisFila}
\noindent\begin{tabularx}{\textwidth}{@{}P{2.8cm} L L C{1.7cm}@{}}
\toprule
\rowcolor{AzulClaro}
\textbf{Alternativa} & \textbf{Efecto en plazo y costo} & \textbf{Efecto en el beneficio} & \textbf{Viable} \\
\midrule
A. Continuar con el alcance completo & Termina el 12/02/2027; costo pronosticado entre S/ 12,537.51 y S/ 15,406.80. & Llega fuera del semestre; depende del convenio. & \etiqueta{SevAlta}{No} \\
B. Reducir a un prototipo con datos simulados & Cabe en el semestre, pero sin piloto ni reportes reales. & Sin reportes ciudadanos no hay beneficio para los usuarios; el VAN económico deja de sustentarse. & \etiqueta{SevAlta}{No} \\
C. Excluir Horus Sentinel & Libera 8 horas semanales por integrante y S/ 807.33 de aporte por integrante. & Se protege el beneficio de CasaScore, cuyo VAN económico es 3.6 veces mayor. & \etiqueta{Rama4}{Sí} \\
\bottomrule
\end{tabularx}
}

\Needspace{10\baselineskip}
\section{Recomendación del Project Manager}
Se recomienda al equipo como sponsor la \textbf{alternativa C: excluir Horus Sentinel} de la cartera, archivar sus entregables de inicio y planificación, y mantener las líneas base de CasaScore. La decisión se somete a votación en la reunión de control del martes 22/09/2026 y, de aprobarse, se registra en el acta de exclusión.

\Needspace{22\baselineskip}
\section{Firmas}
\noindent En señal de conformidad con el presente documento, los integrantes del equipo firman a continuación:

\vspace{1.5cm}
\begin{center}
\begin{tabular}{c c c}
    \parbox{6.5cm}{\centering \rule{6.0cm}{0.6pt}\\[5pt] \textbf{Clyde Neil Paricahua Pari}\\[2pt] {\footnotesize Project Manager}}
    & \hspace{0.8cm} &
    \parbox{6.5cm}{\centering \rule{6.0cm}{0.6pt}\\[5pt] \textbf{Ruth Karina Apaza Solis}\\[2pt] {\footnotesize Analista de Planificación}}
    \\[2.0cm]
    \parbox{6.5cm}{\centering \rule{6.0cm}{0.6pt}\\[5pt] \textbf{Eliseo Tarqui Ajahuana}\\[2pt] {\footnotesize Especialista Técnico / Arquitectura}}
    & \hspace{0.8cm} &
    \parbox{6.5cm}{\centering \rule{6.0cm}{0.6pt}\\[5pt] \textbf{Viviana Mary Quisbert Quispe}\\[2pt] {\footnotesize Gestora de Comunicación, Riesgos y Documentación}}
\end{tabular}
\end{center}

\end{document}
